\chapter{Subcalibrations and minimizing Simons cones}\label{sec:subcalibrations-simons}
The signed-divergence method proves perimeter minimization even when a divergence-free calibration is unavailable. We use outward normals, so $D\chi_E=-\nu_E|D\chi_E|$. The model calculation is \cite{ChMin}, Theorem 17.1 and Exercise 17.1, pp.~48--50. We include perturbations on both sides of the cone and justify the field at its singular point. The comparison class here consists of sets of locally finite perimeter.

\begin{lemma}[Locality of perimeter under compact perturbations]\label{lem:perimeter-perturbation-locality}
Suppose $E,F$ have locally finite perimeter in $U\subset\mathbb R^D$ and $\chi_E=\chi_F$ almost everywhere on $U\setminus K$, where $K\Subset U$ is compact. Then their distributional derivatives and perimeter measures agree on $U\setminus K$. If $K\subset W\Subset U$ is open and $\phi\in C_c^1(W)$ equals one near $K$, then
\[
 P(F;W)-P(E;W)=\int\phi\,d(|D\chi_F|-|D\chi_E|).
\]
\end{lemma}
\begin{proof}
Tests supported outside $K$ have identical integrals against the two characteristic functions, hence identical distributional derivatives there. Uniqueness of the representing vector measure and its total variation gives equality of the perimeter measures on that open set. The signed difference is therefore supported in $K$. Since the measures are finite on $\overline W$, integration against $\phi$ agrees with their total difference on $W$.
\end{proof}

\begin{lemma}[Removing one singular point from a flux identity]\label{lem:singular-field-flux}
Let $D\geq2$, $a\in U$, and let $X$ be bounded and $C^1$ on $U\setminus\{a\}$, with its classical divergence in $L^1_{\mathrm{loc}}(U)$. Give $X(a)$ any bounded value. For every locally finite-perimeter set $F$ and $\phi\in C_c^1(U)$,
\[
 \int_F\bigl(\phi\operatorname{div}X+D\phi\cdot X\bigr)\,dx
 =\int_{\partial^*F}\phi X\cdot\nu_F\,d\mathcal H^{D-1}.
\]
\end{lemma}
\begin{proof}
Let $\eta_\varepsilon$ be zero on $B_\varepsilon(a)$, one outside $B_{2\varepsilon}(a)$, with $|D\eta_\varepsilon|\leq C/\varepsilon$. Apply Theorem \ref{thm:de-giorgi} to the compactly supported $C^1$ field $\eta_\varepsilon\phi X$. The extra bulk term has magnitude at most $C\varepsilon^{-1}\mathcal L^D(B_{2\varepsilon})=O(\varepsilon^{D-1})$, tending to zero. All other bulk terms converge by boundedness and the assumed integrability of the divergence. De Giorgi's representation implies that $|D\chi_F|$ has no point atoms when $D\geq2$. Dominated convergence therefore also applies to the boundary flux. This proves the formula without a point-supported divergence term.
\end{proof}

\begin{definition}[A signed subcalibration for a perimeter set]\label{def:signed-subcalibration}
For a locally finite-perimeter set $E\subset U\subset\mathbb R^D$, $D\geq2$, a signed subcalibration is a field $X$ that is $C^1$ in $U$ or satisfies Lemma \ref{lem:singular-field-flux} at one exceptional point, and obeys
\[
 |X|\leq1,\qquad X\cdot\nu_E=1\quad |D\chi_E|\text{-almost everywhere},
\]
\[
 \operatorname{div}X\leq0\text{ almost everywhere in }E,
 \qquad \operatorname{div}X\geq0\text{ almost everywhere in }U\setminus E.
\]
The boundary condition saturates the flux bound. Unlike a calibration arising from a closed form, the divergence may be nonzero. The sign convention agrees with the outward field in \cite{ChMin}, Section 17, pp.~49--50.
\end{definition}

\begin{proposition}[The exact perimeter-gap identity]\label{prop:subcalibration-gap}
Let $X$ be a signed subcalibration for $E$, and let $F$ have locally finite perimeter with $\chi_F=\chi_E$ almost everywhere outside a compact $K\Subset U$. For every open $K\subset W\Subset U$,
\[
 P(F;W)-P(E;W)
 =\int_W(\chi_F-\chi_E)\operatorname{div}X\,dx
 +\int_{\partial^*F\cap W}(1-X\cdot\nu_F)\,d\mathcal H^{D-1}.
\]
Both terms on the right are nonnegative.
\end{proposition}
\begin{proof}
Choose $0\leq\phi\leq1$ compactly supported in $W$ and equal to one near $K$. Subtract the flux identities for $F$ and $E$, using Lemma \ref{lem:singular-field-flux} if necessary. The term $(\chi_F-\chi_E)D\phi\cdot X$ is zero. The difference of boundary flux measures is supported in $K$, by the equality of vector derivatives in Lemma \ref{lem:perimeter-perturbation-locality}. Thus the cutoffs can be removed from the difference, giving
\[
 \int_{\partial^*F\cap W}X\cdot\nu_F\,d\mathcal H^{D-1}
 -\int_{\partial^*E\cap W}X\cdot\nu_E\,d\mathcal H^{D-1}
 =\int_W(\chi_F-\chi_E)\operatorname{div}X.
\]
The second boundary integral equals $P(E;W)$ by the saturated flux. Subtract the first from $P(F;W)$ to obtain the identity. The bulk term is nonnegative separately on $F\setminus E$ and $E\setminus F$ by the divergence signs. The boundary term is nonnegative because $|X|\leq1$.
\end{proof}

\begin{corollary}[Perimeter minimality and strictness from a subcalibration]\label{cor:subcalibration-minimizes}
A set admitting a signed subcalibration minimizes perimeter under every compactly supported perturbation by a set of locally finite perimeter. If in addition $\operatorname{div}X<0$ almost everywhere in $E$ and $\operatorname{div}X>0$ almost everywhere outside $E$, then
\[
 P(F;W)>P(E;W)
\]
for every such competitor with $\mathcal L^D(E\mathbin{\triangle}F)>0$, in each relatively compact open $W$ containing its compact perturbation support.
\end{corollary}
\begin{proof}
Use Proposition \ref{prop:subcalibration-gap}. Under the strict signs, its bulk integrand is strictly positive almost everywhere on the symmetric difference and zero elsewhere. Its integral is positive whenever that difference has positive measure. Equality almost everywhere is the appropriate exception because perimeter depends only on the characteristic function as an $L^1_{\mathrm{loc}}$ class. No nesting assumption between $E$ and $F$ is used.
\end{proof}

\begin{proposition}[Perimeter and normal of the equal-factor cone]\label{prop:simons-cone-perimeter}
For $k\geq2$, set
\[
 E_k=\{(x,y)\in\mathbb R^k\times\mathbb R^k:|x|<|y|\},
 \qquad C_k=\{|x|=|y|\}.
\]
The set $E_k$ has locally finite perimeter, and
\[
 |D\chi_{E_k}|=\mathcal H^{2k-1}\llcorner(C_k\setminus\{0\}).
\]
Away from zero its outward unit normal is
\[
 \nu_{E_k}(x,y)=\frac{(x,-y)}{\sqrt{|x|^2+|y|^2}}.
\]
Its link is $L=\mathbb S^{k-1}(1/\sqrt2)\times\mathbb S^{k-1}(1/\sqrt2)$, and
\[
 P(E_k;B_R(0))=\frac{\mathcal H^{2k-2}(L)}{2k-1}R^{2k-1}.
\]
\end{proposition}
\begin{proof}
The function $|x|^2-|y|^2$ has nonzero gradient on $C_k\setminus\{0\}$, giving a smooth boundary there and the displayed outward normal. The cone parametrization $F(r,p)=rp$ has area element $r^{2k-2}dr\,dA_L$ by Proposition \ref{prop:cone-geometry}, giving the area formula and local finiteness away from zero.

To identify the full distributional derivative, apply the smooth divergence theorem with a cutoff vanishing on $B_\varepsilon(0)$. For a fixed compactly supported test field, the term from the cutoff gradient is $O(\varepsilon^{2k-1})$ by its boundedness and the ambient ball volume. The cone boundary integrals converge because the computed area near zero tends to zero. Thus $\int_{E_k}\operatorname{div}Y=\int_{C_k\setminus\{0\}}Y\cdot\nu_{E_k}\,d\mathcal H^{2k-1}$. This gives the derivative measure with no atom at the vertex; its variation is the claimed area measure since the normal has unit length.
\end{proof}

\begin{proposition}[The Simons subcalibration and its divergence]\label{prop:simons-subcalibration-divergence}
Write $r=|x|$, $s=|y|$, and define away from zero
\[
 X(x,y)=\frac{(r^2x,-s^2y)}{(r^6+s^6)^{1/2}}.
\]
Then $|X|=1$, $X=\nu_{E_k}$ on $C_k\setminus\{0\}$, and
\[
 \operatorname{div}X=
 \frac{(r^4-s^4)\bigl((k-1)(r^4+s^4)-(k+2)r^2s^2\bigr)}{(r^6+s^6)^{3/2}}.
\]
For $k\geq4$, the divergence has the strict sign of $r-s$ off $C_k$, and $|\operatorname{div}X|\leq A_k/|(x,y)|$ away from zero for a finite constant $A_k$.
\end{proposition}
\begin{proof}
The numerator has squared length $r^6+s^6$, and on the equal-radius cone its normalization is the normal in Proposition \ref{prop:simons-cone-perimeter}. Put $A=r^6+s^6$. Differentiating the $x$ block and summing over its $k$ coordinates gives
\[
 \operatorname{div}_x(r^2xA^{-1/2})
 =(k+2)r^2A^{-1/2}-3r^8A^{-3/2}
 =\frac{(k-1)r^8+(k+2)r^2s^6}{A^{3/2}}.
\]
Subtract the analogous $y$ expression and factor $r^4-s^4$ to obtain the formula. The remaining bracket can be rewritten as
\[
 (k-1)(r^2-s^2)^2+(k-4)r^2s^2.
\]
It is positive whenever $r\ne s$ and $k\geq4$, including $k=4$. Finally $X$ is smooth on the unit sphere and homogeneous of degree zero, so its divergence is bounded there and homogeneous of degree minus one. This proves the integrable bound near the vertex. This calculation expands the field in \cite{ChMin}, pp.~49--50, using block dimension $k$ rather than the source's $n+1$.
\end{proof}

\begin{theorem}[Strict perimeter minimization by the equal-factor Simons cone]\label{thm:simons-cone-minimizing}
Let $k\geq4$ and $E_k=\{|x|<|y|\}\subset\mathbb R^{2k}$. For every set $F$ of locally finite perimeter such that $\chi_F=\chi_{E_k}$ almost everywhere outside a compact set $K$, and every relatively compact open $W$ containing $K$,
\[
 P(E_k;W)\leq P(F;W).
\]
The inequality is strict if $\mathcal L^{2k}(E_k\mathbin{\triangle}F)>0$. This is the finite-perimeter-set version of \cite{ChMin}, Theorem 17.1 and Exercise 17.1, pp.~48--50, extending the source's worked one-sided comparison to arbitrary compact perturbations.
\end{theorem}
\begin{proof}
Proposition \ref{prop:simons-cone-perimeter} supplies locally finite perimeter and its normal. The field in Proposition \ref{prop:simons-subcalibration-divergence} saturates that normal flux, has norm one, and has integrable divergence near its only singular point. Set its value at the origin to zero. Lemma \ref{lem:singular-field-flux} therefore justifies its use for every competitor $F$. The divergence is negative in $E_k$ and positive in its complement except on $C_k$, a Lebesgue-null set. It is a signed subcalibration with the strict sign conditions. Corollary \ref{cor:subcalibration-minimizes} gives both claims. This comparison class is that of finite-perimeter sets; an extension to all integral currents with arbitrary multiplicity would require an additional decomposition argument.
\end{proof}

\begin{corollary}[A singular perimeter minimizer in ambient dimension eight]\label{cor:simons-eight-dimensional-example}
The boundary $C_4\subset\mathbb R^8$ is a seven-dimensional locally perimeter-minimizing cone, smooth away from zero and singular at zero. Thus perimeter minimization does not imply everywhere smooth boundary in this dimension.
\end{corollary}
\begin{proof}
Theorem \ref{thm:simons-cone-minimizing} gives minimization, and Proposition \ref{prop:simons-cone-perimeter} gives smoothness away from the vertex. If the cone were a smooth hypersurface at zero, its rescalings would converge to its tangent plane. Every rescaling is the same cone, so it would be a plane. It is not: for nonzero $u\in\mathbb R^4$, both $(u,u)$ and $(u,-u)$ belong to it, whereas their sum $(2u,0)$ does not. This proves singularity. The conclusion does not classify other minimizing cones or assert a boundary regularity theorem in lower dimensions.
\end{proof}

\begin{corollary}[Stationarity and punctured stability of the minimizing cone]\label{cor:simons-minimizer-stationary}
For $k\geq4$, the multiplicity-one varifold $|C_k|$ is stationary on all of $\mathbb R^{2k}$, including at zero, and its smooth part is stable for compactly supported normal variations away from zero.
\end{corollary}
\begin{proof}
The flow of a smooth compactly supported ambient field sends $E_k$ to a finite-perimeter set equal to $E_k$ outside a compact set. Its boundary measure is the image cone area, as follows from the area formula and the zero mass of the single vertex. Theorem \ref{thm:simons-cone-minimizing} says that the local perimeter has a minimum at time zero. Differentiating its tangential Jacobian under the locally finite cone measure gives zero first variation for every such field; bounded flow derivatives on the compact support justify the differentiation. Thus stationarity holds even for fields nonzero at the vertex.

A smooth compactly supported normal variation on the punctured cone extends to an ambient field in a tubular neighborhood away from zero. The second derivative of area at its minimum is nonnegative. The usual smooth second-variation formula gives the stability inequality in Definition \ref{def:cone-punctured-stability}. Compare the independent spectral calculation in Example \ref{ex:clifford-cone-stability}; this argument uses the stronger minimizing property for equal factors.
\end{proof}
